\documentclass[a4paper,11pt]{article}

\usepackage[margin=2.4cm]{geometry}
\usepackage{amsmath,amssymb,bm}
\usepackage[dvipsnames]{xcolor}
\usepackage{enumitem}
\usepackage{xr}
\externaldocument{EFT_lectures}   % equation numbers refer to the lecture notes
\usepackage[colorlinks=true,linkcolor=RoyalBlue,citecolor=ForestGreen]{hyperref}

\newcommand{\vx}{\bm{x}}
\newcommand{\vq}{\bm{q}}
\newcommand{\vk}{\bm{k}}
\newcommand{\vu}{\bm{u}}
\newcommand{\vPsi}{\bm{\Psi}}
\newcommand{\nhat}{\hat{\bm{z}}}
\newcommand{\khat}{\hat{\bm{k}}}
\def\cH{\mathcal{H}}
\def\Plin{P_{\rm L}}
\def\knl{k_{\rm NL}}
\newcommand{\dd}{\mathrm{d}}
\def\dirac{\delta_{\rm D}}
\def\Om{\Omega_m}
\def\G{\mathcal{G}}
\def\lcdm{$\Lambda$CDM}
\def\vsv{\bm{s}}

\newenvironment{solution}[1]
  {\par\medskip\noindent\textbf{#1}\par\nobreak\smallskip}
  {\par\medskip}

\title{Cosmological Perturbation Theory from Matter to Galaxies\\[2mm]
\large Exercise solutions}
\author{Nhat-Minh Nguyen}
\date{\today}

\begin{document}
\maketitle

\noindent Solutions to the exercises in the lecture notes. Equation numbers
refer to the notes; equations numbered here are local to each solution.

%=====================================================================
\section*{Lecture 1}

\begin{solution}{Exercise 1.1 (the Newtonian equation of motion).}
Conformal time is not an affine parameter, so we cannot use the affine form of
the geodesic equation directly. Parametrizing the worldline by $x^0 = \tau$ and
eliminating the affine parameter gives the standard result
\begin{equation*}
\ddot x^i + \Gamma^i{}_{\alpha\beta}\dot x^\alpha\dot x^\beta
- \dot x^i\,\Gamma^0{}_{\alpha\beta}\dot x^\alpha\dot x^\beta = 0 ,
\end{equation*}
with overdots denoting $\dd/\dd\tau$ and $\dot x^0 = 1$. The last term is the
one that is easy to forget; dropping it doubles the Hubble drag. For the metric
$\dd s^2 = a^2[-(1+2\Phi)\dd\tau^2 + (1-2\Phi)\dd\vx^2]$, to first order in
$\Phi$ and leading order in $u \ll 1$, the relevant symbols are
\begin{equation*}
\Gamma^i{}_{00} = \partial_i\Phi ,
\qquad
\Gamma^i{}_{0j} = \left(\cH - \dot\Phi\right)\delta^i{}_j ,
\qquad
\Gamma^0{}_{00} = \cH + \dot\Phi .
\end{equation*}
Substituting,
\begin{equation*}
\ddot x^i + 2\left(\cH-\dot\Phi\right)\dot x^i + \partial_i\Phi
- \left(\cH+\dot\Phi\right)\dot x^i = 0
\qquad\Longrightarrow\qquad
\ddot{\vx} + \cH\,\dot{\vx} = -\nabla\Phi + 3\dot\Phi\,\dot{\vx} .
\end{equation*}
The final term carries one power of the potential and one of the velocity, so it
is higher order and we drop it, leaving
\begin{equation*}
\ddot{\vx} + \cH\,\dot{\vx} = -\nabla\Phi ,
\end{equation*}
the Newtonian equation for a particle in an expanding background, with $\cH\dot{\vx}$
the Hubble drag. A quick check on the drag coefficient: with $\Phi = 0$ the
solution is $\dot{\vx} \propto 1/a$, the familiar decay of peculiar velocities.
Nothing here required $k \gg \cH$; that limit was used earlier to reduce the
Einstein constraint to the Poisson equation, which is what makes $\Phi$ the
Newtonian potential.
\end{solution}

\begin{solution}{Exercise 1.2 (the sub-horizon limit).}
\emph{(i)} The two groups in \eqref{eq:energy_constraint} are $k^2\phi$ and
$3\cH(\dot\phi + \cH\psi)$. With $\phi=\psi=\Phi$ and $\dot\Phi \lesssim \cH\Phi$,
the second is bounded by $3\cH(\cH\Phi + \cH\Phi) = 6\cH^2\Phi$, so
\begin{equation*}
\frac{\text{discarded}}{\text{retained}} \;\lesssim\; 6\left(\frac{\cH}{k}\right)^2 .
\end{equation*}
For the growing mode in matter domination the bound is saturated from below
rather than above: $\Phi$ is constant, so $\dot\Phi = 0$ exactly and the ratio is
$3(\cH/k)^2$. Either way the suppression is two powers of $\cH/k$, and the limit
$k \gg \cH$ leaves $-k^2\Phi = 4\pi Ga^2\bar\rho\,\delta$.

\emph{(ii)} $(\cH_0/k)^2 = (3.3\times10^{-4}/k)^2$ gives $1.1\times10^{-5}$ at
$k = 0.1\,h\,{\rm Mpc}^{-1}$ and $1.1\times10^{-1}$ at
$k = 10^{-3}\,h\,{\rm Mpc}^{-1}$. So on the scales this course actually uses the
Newtonian field equation is good to a part in $10^5$, while at
$k \sim 10^{-3}$ the neglected term is a ten-percent effect. That is the
low-$k$ edge of the window in the Comment, and this is what closes it.

\emph{(iii)} Because a bound is a claim about the solution, while an identity is
a claim about the equations. The bound $\dot\Phi \lesssim \cH\Phi$ holds for the
growing mode and fails for the decaying one, and one has to check it again in
every new regime --- during a rapid transition, for a modified-gravity model, for
an isocurvature initial condition. The algebraic form has no such exposure: the
time derivative has been eliminated using a second Einstein equation, so the
$\mathcal{O}(\cH/k)$ suppression is visible in the equation itself and holds
wherever the Einstein equations do. Removing an assumption is worth more than
verifying it, because verification does not travel.
\end{solution}

\begin{solution}{Exercise 1.3 ($D_2$ in EdS).}
In EdS, $a \propto \tau^2$, $\cH = 2/\tau$, $\Om = 1$, and $D_+ = a$. Try
$D_2 = C a^2 \propto \tau^4$. Then $\dot D_2 = 4C\tau^3$ and
$\ddot D_2 = 12C\tau^2$, while $D_+^2 = a^2 \propto \tau^4$. Normalising
$a = \tau^2$, equation \eqref{eq:D2} reads
\begin{equation*}
12C\tau^2 + \frac{2}{\tau}\,4C\tau^3 - \frac{3}{2}\cdot\frac{4}{\tau^2}\,C\tau^4
= -\frac{3}{2}\cdot\frac{4}{\tau^2}\,\tau^4 ,
\end{equation*}
that is $12C + 8C - 6C = -6$, so $14C = -6$ and $C = -3/7$. Hence
$D_2 = -\tfrac37 D_+^2$, the EdS limit quoted in \eqref{eq:D2}. Note that the
associated growth rate must be defined as $f_2 = \dd\ln|D_2|/\dd\ln a$, cf.\
\eqref{eq:f2_def}: with this sign convention $D_2<0$ and $\ln D_2$ does not
exist. In EdS, $|D_2| \propto a^2$ gives $f_2 = 2$.

The negative sign is a bookkeeping artifact of the convention
$\vPsi^{(2)} = +\nabla_q\varphi^{(2)}$, not a statement that second order opposes
first. Check with a spherical patch, where $\varphi^{(1)}_{,ij} = (\delta_1/3)\delta_{ij}$
and the tidal invariant is $\delta_1^2/3$, so \eqref{eq:2lpt_potential} gives
$\nabla_q\cdot\vPsi^{(2)} = \nabla^2_q\varphi^{(2)} = -\delta_1^2/7$. This has the
\emph{same} sign as $\nabla_q\cdot\vPsi^{(1)} = -\delta_1$: the second-order
displacement reinforces the infall. Equivalently, expanding the spherical
Zel'dovich density $(1-D_+\delta_1/3)^{-3}$ gives $\delta^{(2)} = \tfrac23\delta_1^2$,
whereas the correct second-order result --- the angular average of the
second-order Eulerian kernel $F_2$ of \eqref{eq:F2G2}, which you meet in
Lecture~2 --- is $\tfrac{17}{21}\delta_1^2$. Since $17/21 > 2/3$, adding the tidal term makes
collapse \emph{faster} than the straight-line Zel'dovich estimate, not slower.
\end{solution}

\begin{solution}{Exercise 1.4 (the Zel'dovich approximation is exact in 1D).}
In one dimension, Gauss's law says the gravitational field at a point depends
only on the total mass on either side of it. Consider a sheet with Lagrangian
coordinate $q$ displaced to $x = q + \Psi(q,\tau)$. Before shell crossing the
ordering of sheets is preserved, so the mass on either side of our sheet is the
same as it was initially: it does not depend on how far anything has moved.

The peculiar force is therefore fixed by the initial configuration alone. Make
this explicit by integrating the 1D Poisson equation
$\partial_x^2\Phi = \tfrac32\Om\cH^2\delta$ once. Using
$1 + \delta = (1 + \partial_q\Psi)^{-1}$ and $\dd x = (1+\partial_q\Psi)\dd q$,
the mass integral collapses exactly:
\begin{equation*}
\int \delta\,\dd x
= \int \left[\frac{1}{1+\partial_q\Psi} - 1\right](1+\partial_q\Psi)\,\dd q
= -\int \partial_q\Psi\,\dd q = -\Psi ,
\end{equation*}
measuring $\Psi$ from a point where it vanishes. Note that every power of
$\partial_q\Psi$ cancels: this is the algebraic content of the statement that the
enclosed mass does not change. Hence
\begin{equation*}
\partial_x\Phi = -\frac{3}{2}\Om\cH^2\,\Psi(q,\tau) ,
\end{equation*}
with no higher-order corrections. The sign is the physically sensible one: a
sheet displaced to $\Psi < 0$, having fallen inward, feels an inward force
$-\partial_x\Phi < 0$. Substituting into the exact equation of motion
\eqref{eq:newtonian_eom} gives
\begin{equation*}
\ddot\Psi + \cH\dot\Psi - \frac{3}{2}\Om\cH^2\Psi = 0 ,
\end{equation*}
which is exactly the linear growth equation \eqref{eq:growth}. So
$\Psi = D_+(\tau)\,\Psi_0(q)$ solves the full nonlinear problem, and the
Zel'dovich displacement is exact until trajectories first cross.

In three dimensions this fails because the force depends on the full
three-dimensional mass distribution, which the displacement itself changes: the
tidal term of \eqref{eq:2lpt_potential} is precisely the leading correction.
\end{solution}

%=====================================================================
\section*{Lecture 2}

\begin{solution}{Exercise 2.1 ($F_2$ from the recursion).}
At $n = 2$ only $m = 1$ contributes to \eqref{eq:Fn_recursion}, and
$F_1 = G_1 = 1$. With $(2n+3)(n-1) = 7$ and $2n+1 = 5$,
\begin{equation*}
F_2(\vq_1,\vq_2) = \frac{5\,\alpha(\vq_1,\vq_2) + 2\,\beta(\vq_1,\vq_2)}{7},
\qquad
G_2(\vq_1,\vq_2) = \frac{3\,\alpha(\vq_1,\vq_2) + 4\,\beta(\vq_1,\vq_2)}{7}.
\end{equation*}
Now symmetrize. Writing $\mu \equiv \hat{\vq}_1\cdot\hat{\vq}_2$,
\begin{equation*}
\frac{\alpha(\vq_1,\vq_2)+\alpha(\vq_2,\vq_1)}{2}
= 1 + \frac{\mu}{2}\left(\frac{q_1}{q_2}+\frac{q_2}{q_1}\right),
\qquad
\beta(\vq_1,\vq_2) = \frac{\mu}{2}\left(\frac{q_1}{q_2}+\frac{q_2}{q_1}\right)
+ \mu^2 ,
\end{equation*}
where the second follows from
$|\vq_1+\vq_2|^2 = q_1^2+q_2^2+2q_1q_2\mu$. Substituting,
\begin{equation*}
F_2^{\rm s} = \frac57 + \frac{\mu}{2}\left(\frac{q_1}{q_2}+\frac{q_2}{q_1}\right)
+ \frac{2}{7}\mu^2 ,
\end{equation*}
and likewise $G_2^{\rm s}$ with $5/7 \to 3/7$ and $2/7 \to 4/7$, which is
\eqref{eq:F2G2}.

For the soft limit set $\vq_1 = \vq \to 0$ and $\vq_2 = \vk-\vq \to \vk$. The
dipole term dominates because $q_2/q_1 \to \infty$:
\begin{equation*}
F_2^{\rm s}(\vq,\vk-\vq) \simeq \frac{\mu}{2}\frac{q_2}{q_1}
= \frac{1}{2}\frac{\vq\cdot\vk}{qk}\cdot\frac{k}{q}
= \frac{\vk\cdot\vq}{2q^2} ,
\end{equation*}
which is kernel property 3. It diverges as $1/q$: a long-wavelength mode moves
short-wavelength structure bodily, and the displacement grows as the mode gets
longer.
\end{solution}

\begin{solution}{Exercise 2.2 (the infrared cancellation).}
\emph{The $P_{22}$ side.} Insert the soft limit of Exercise 2.1 into
\eqref{eq:P22}. For $q \ll k$ the second spectrum is
$\Plin(|\vk-\vq|) \to \Plin(k)$, so
\begin{equation*}
P_{22}(k) \supset 2\int_{\vq} \frac{(\vk\cdot\vq)^2}{4q^4}\,
\Plin(q)\,\Plin(k)
= \frac{\Plin(k)}{2}\int\!\frac{\dd^3q}{(2\pi)^3}\,
\frac{k^2q^2\mu_q^2}{q^4}\,\Plin(q) ,
\end{equation*}
with $\mu_q$ the angle between $\vq$ and $\vk$. Using
$\int \dd\Omega\,\mu_q^2 = 4\pi/3$,
\begin{equation*}
P_{22}(k) \supset \frac{\Plin(k)k^2}{2}\cdot\frac{1}{3}\cdot
\frac{1}{2\pi^2}\int\!\dd q\,\Plin(q)
= \frac{1}{2}\,k^2\sigma_v^2\,\Plin(k) ,
\end{equation*}
with $\sigma_v^2$ as defined in \eqref{eq:P13_UV}. The same contribution comes
from the mirror region $|\vk-\vq| \to 0$, where the roles of the two spectra are
exchanged, so the total soft contribution is
$+\,k^2\sigma_v^2\,\Plin(k)$.

\emph{The $P_{13}$ side.} The angle-averaged kernel in the same limit is
$\langle F_3(\vk,\vq,-\vq)\rangle \to -k^2/(18q^2)$, so from \eqref{eq:P13}
\begin{equation*}
P_{13}(k) \supset 6\,\Plin(k)\,\frac{1}{2\pi^2}\int\!\dd q\;q^2\,
\left(-\frac{k^2}{18q^2}\right)\Plin(q)
= -\,\frac{k^2}{6\pi^2}\,\Plin(k)\int\!\dd q\,\Plin(q)
= -\,k^2\sigma_v^2\,\Plin(k) ,
\end{equation*}
where the factor $1/2\pi^2$ comes from $\int\dd^3q/(2\pi)^3$ after the angular
average. The two contributions cancel exactly in $P_{22}+P_{13}$.

The physical statement is that $\sigma_v^2$ measures a uniform displacement of
everything, which cannot affect a correlation function evaluated at a single
time. This is the equivalence principle, and the cancellation it enforces is
exact for any shape of $\Plin$ --- including one with BAO.

What it removes is the strictly uniform piece. A feature of intrinsic scale
$k_{\rm osc}$ supplies a second scale, and a mode with $q \gtrsim k_{\rm osc}$
can be soft relative to $k$ while $q/k_{\rm osc}$ is of order unity. Expanding
in $q/k_{\rm osc}$ then does not terminate: an infinite set of finite
IR-enhanced terms survives the cancellation, and resumming them is what damps
the acoustic peak (Lecture 3). Note the logic --- the BAO feature does not
violate the cancellation; it makes the \emph{finite} remainder large.
\end{solution}

\begin{solution}{Exercise 2.3 (LPT reproduces $F_2$).}
Work in EdS and drop the growth factors, i.e. set $D_+ = 1$ and
$D_2 = -3/7$. In Fourier space the first-order displacement is
$\vPsi^{(1)}(\vk) = (i\vk/k^2)\,\delta^{(1)}(\vk)$, so for two modes
$\vq_1, \vq_2$ summing to $\vk$,
\begin{equation*}
\Psi^{(1)}_{k,k}(\vq_i) \to -\delta^{(1)}(\vq_i) ,
\qquad
\Psi^{(1)}_{i,j}(\vq) \to -\frac{q_iq_j}{q^2}\,\delta^{(1)}(\vq) .
\end{equation*}
Take the four terms of \eqref{eq:delta2_lpt} in turn, symmetrizing in
$\vq_1 \leftrightarrow \vq_2$ and writing $\mu = \hat{\vq}_1\cdot\hat{\vq}_2$.

\emph{(i) $\tfrac12(\Psi^{(1)}_{k,k})^2$} contributes $\tfrac12$, a pure
constant.

\emph{(ii) $\tfrac12\Psi^{(1)}_{i,j}\Psi^{(1)}_{j,i}$} contributes
$\tfrac12(\hat{\vq}_1\cdot\hat{\vq}_2)^2 = \tfrac12\mu^2$.

\emph{(iii) $-\Psi^{(2)}_{k,k}$.} Since $\vPsi^{(2)} = \nabla_q\varphi^{(2)}$,
we need $\Psi^{(2)}_{k,k} = \nabla_q^2\varphi^{(2)}$, which
\eqref{eq:2lpt_potential} gives as $D_2/D_+^2$ times the tidal invariant. Write
that invariant as
$\sum_{i>j}[\varphi_{,ii}\varphi_{,jj} - \varphi_{,ij}^2]
= \tfrac12[(\nabla^2\varphi)^2 - \varphi_{,ij}\varphi_{,ij}]$. The first square
has Fourier kernel $1$ and the second $\mu^2$, so the invariant has kernel
$\tfrac12(1-\mu^2)$. With $D_2/D_+^2 = -3/7$,
\begin{equation*}
-\Psi^{(2)}_{k,k} \;\longrightarrow\; \frac{3}{7}\cdot\frac12\left(1-\mu^2\right)
= \frac{3}{14}\left(1-\mu^2\right).
\end{equation*}
Adding (i), (ii) and (iii),
\begin{equation*}
\frac12 + \frac12\mu^2 + \frac{3}{14}\left(1-\mu^2\right)
= \frac{7+3}{14} + \frac{7-3}{14}\,\mu^2
= \frac{5}{7} + \frac{2}{7}\mu^2 ,
\end{equation*}
which is the monopole and quadrupole of $F_2$.

\emph{(iv) The transport term $\Psi^{(1)}_j\partial_j\Psi^{(1)}_{k,k}$} is the
one that supplies the dipole. In Fourier space
$\Psi^{(1)}_j(\vq_1)\,(iq_{2j})\,(-\delta^{(1)}(\vq_2))$ gives, after
symmetrizing,
\begin{equation*}
\frac{1}{2}\left(\frac{\vq_1\cdot\vq_2}{q_1^2} + \frac{\vq_1\cdot\vq_2}{q_2^2}\right)
= \frac{\mu}{2}\left(\frac{q_2}{q_1}+\frac{q_1}{q_2}\right).
\end{equation*}

Adding the four pieces reproduces $F_2$ of \eqref{eq:F2G2} exactly. The lesson
is worth stating: the constant and $\mu^2$ parts come from the volume
distortion encoded in the Jacobian, while the dipole comes entirely from
evaluating the field at the displaced position. That dipole is the advection
term whose infrared behaviour drove Exercise 2.2.

If the bookkeeping above feels fragile, the cleanest route is the exact map
\begin{equation*}
\delta(\vk) = \int\!\dd^3q\; e^{-i\vk\cdot\vq}
\left(e^{-i\vk\cdot\vPsi(\vq)}-1\right),
\end{equation*}
which follows from \eqref{eq:jacobian} by Fourier transforming
$1+\delta$. Expanding the exponential to second order in $\vPsi$ and inserting
$\vPsi^{(1)}$, $\vPsi^{(2)}$ produces $F_2$ with no separate transport term,
because the Eulerian position is built in from the start.
\end{solution}

\begin{solution}{Exercise 2.4 (the running of $c_s^2$).}
Collect the terms in \eqref{eq:renormalized} that carry the shape $k^2\Plin(k)$.
$P_{22}$ has none: its leading UV behaviour goes as $k^4$ with no factor of
$\Plin(k)$, because momentum conservation forces $F_2 \propto k^2$ in the
squeezed limit and the kernel enters squared. So at this order only two terms
depend on $\Lambda$, the UV limit \eqref{eq:P13_UV} of the loop and the
counterterm:
\begin{equation*}
P_{\rm 1\text{-}loop}\Big|_{k^2\Plin}
= -\frac{61}{105}\,k^2\sigma_v^2(\Lambda)\,\Plin(k)
\;-\; 2\,c_s^2(\Lambda)\,k^2\,\Plin(k) .
\end{equation*}
Demanding $\partial P_{\rm 1\text{-}loop}/\partial\ln\Lambda = 0$ and dividing
out the common $k^2\Plin(k)$,
\begin{equation*}
-\frac{61}{105}\,\frac{\dd\sigma_v^2}{\dd\ln\Lambda}
- 2\,\frac{\dd c_s^2}{\dd\ln\Lambda} = 0
\qquad\Longrightarrow\qquad
\frac{\dd c_s^2}{\dd\ln\Lambda} = -\frac{61}{210}\,\frac{\dd\sigma_v^2}{\dd\ln\Lambda} .
\end{equation*}
With $\sigma_v^2(\Lambda) = \tfrac{1}{6\pi^2}\int_0^\Lambda\!\dd q\,\Plin(q)$
the derivative is $\dd\sigma_v^2/\dd\ln\Lambda = \Lambda\Plin(\Lambda)/6\pi^2$,
which is the quoted running. Integrating,
\begin{equation*}
c_s^2(\Lambda) = c_{s,\rm ren}^2 - \frac{61}{210}\,\sigma_v^2(\Lambda) ,
\end{equation*}
with $c_{s,\rm ren}^2$ an integration constant. Substituting back,
\begin{equation*}
P_{13} + P_{\rm ctr}
= -\frac{61}{105}k^2\sigma_v^2\Plin
- 2k^2\Plin\left[c_{s,\rm ren}^2 - \frac{61}{210}\sigma_v^2\right]
= -2\,c_{s,\rm ren}^2\,k^2\,\Plin(k) ,
\end{equation*}
and every trace of $\Lambda$ has cancelled.

What is predicted and what is not: the \emph{shape} $k^2\Plin$ and the
\emph{coefficient of the running}, $61/210$, both follow from $F_3$ and are
genuine predictions of the theory. The finite constant $c_{s,\rm ren}^2$ is not.
It is an integration constant, and no manipulation of the dust equations can
produce it, because it encodes the response of scales where those equations do
not hold. Renormalization predicts the cutoff dependence, never the value. That
one number has to be measured --- from simulations, or from the data alongside
the cosmology.
\end{solution}

\begin{solution}{Exercise 2.5 (matter stochasticity starts at $k^4$).}
Expand the Fourier transform of the stochastic density at small $k$:
\begin{equation*}
\delta_\varepsilon(\vk) = \int\!\dd^3x\;e^{-i\vk\cdot\vx}\,\delta_\varepsilon(\vx)
= \int\!\dd^3x\,\delta_\varepsilon
- i k_i\!\int\!\dd^3x\;x_i\,\delta_\varepsilon
- \frac{k_ik_j}{2}\!\int\!\dd^3x\;x_ix_j\,\delta_\varepsilon
+ \ldots
\end{equation*}
Mass conservation kills the first integral and momentum conservation the second,
so the series starts at $\mathcal{O}(k^2)$. The auto-power is quadratic in the
field, hence $P_\varepsilon(k) = \mathcal{O}(k^4)$. Physically: a small-scale
rearrangement of matter that neither creates mass nor shifts the centre of mass
of a patch cannot imprint anything on the patch's monopole or dipole, and the
leading residual costs two powers of $k$. (The surviving second
moment has an unconstrained trace, so that residual is generically an isotropic
$k^2$ piece plus a quadrupole; only the $k^4$ scaling of the power matters
here.)

For galaxies the first constraint is simply false. Galaxy number is not
conserved --- galaxies form, merge, and drop out of a sample --- so
$\int\!\dd^3x\,\delta_\varepsilon \ne 0$, nothing forbids a $k^0$ term, and the
stochastic spectrum starts at a constant. That constant is shot noise. This is
the sense in which shot noise is not a nuisance imported from Poisson counting
but the leading term of an expansion whose matter counterpart happens to be
forbidden.
\end{solution}

%=====================================================================
\section*{Lecture 3}

\begin{solution}{Exercise 3.1 (the Kaiser formula).}
Expand \eqref{eq:rsd_exact} to linear order. The exponential contributes
$-i(k\mu/\cH)u_z(\vx)$, and the factor $[1+\delta_g]$ can be set to unity since
the velocity term is already first order:
\begin{equation*}
\delta_g^s(\vk) = \delta_g(\vk)
- \frac{ik\mu}{\cH}\int\!\dd^3x\;e^{-i\vk\cdot\vx}u_z(\vx)
= \delta_g(\vk) - \frac{ik\mu}{\cH}\,u_z(\vk) .
\end{equation*}
For an irrotational flow $\vu(\vk) = -i\hat{\vk}\,\theta(\vk)/k$, so the
line-of-sight component is $u_z(\vk) = -i\mu\,\theta(\vk)/k$. Then
\begin{equation*}
-\frac{ik\mu}{\cH}\cdot\left(-\frac{i\mu\theta}{k}\right)
= -\frac{\mu^2}{\cH}\,\theta(\vk) .
\end{equation*}
Using $\theta = -f\cH\delta$ from \eqref{eq:growth_rate} this becomes
$+f\mu^2\delta$, and with $\delta_g = b_1\delta$ at linear order,
\begin{equation*}
\delta_g^s(\vk) = \left(b_1 + f\mu^2\right)\delta^{(1)}(\vk) ,
\end{equation*}
which squares to \eqref{eq:kaiser}. The sign is worth checking carefully: it is
positive, so infall enhances clustering along the line of sight. This is what
makes the quadrupole sensitive to the growth rate --- though what it actually
measures is $\beta = f/b_1$, and with the amplitude restored, $f\sigma_8$; see
the solution to Exercise 3.3.
\end{solution}

\begin{solution}{Exercise 3.2 (deriving $Z_2$).}
\emph{(i)} Setting $f = 0$ in \eqref{eq:Z2} kills the $f\mu^2G_2$ term and the
whole second line, leaving
\begin{equation*}
Z_2 \to b_1F_2(\vk_1,\vk_2) + \frac{b_2}{2}
+ b_{\G_2}\left[(\khat_1\cdot\khat_2)^2-1\right] ,
\end{equation*}
the real-space galaxy kernel: gravitational evolution of the linear density
weighted by $b_1$, plus the quadratic bias operators. Redshift space adds no
new bias, only new kinematics.

\emph{(ii)} Now do the algebra. Expand the exponential in \eqref{eq:rsd_exact},
writing $A \equiv (k\mu/\cH)u_z$, so that $e^{-iA}-1 = -iA - \tfrac12 A^2 +
\ldots$, and collect everything second order in $\delta^{(1)}$:
\begin{equation*}
\delta_g^{s(2)}
= \underbrace{\delta_g^{(2)}}_{\rm (a)}
\;\underbrace{-\,\frac{ik\mu}{\cH}\,u_z^{(2)}}_{\rm (b)}
\;\underbrace{-\,\frac{ik\mu}{\cH}\left[u_z\,\delta_g\right]^{(1\times1)}}_{\rm (c)}
\;\underbrace{-\,\frac12\left(\frac{k\mu}{\cH}\right)^{\!2}
\left[u_z u_z\right]^{(1\times1)}}_{\rm (d)} .
\end{equation*}
The one ingredient needed throughout is the linear line-of-sight velocity in
Fourier space. From $u_z(\vk) = -i\mu\,\theta(\vk)/k$ and
$\theta = -\cH f\delta$,
\begin{equation*}
u_z^{(1)}(\vk_i) = -\frac{i\mu_i}{k_i}\left(-\cH f\,\delta^{(1)}_i\right)
= \frac{i\cH f\mu_i}{k_i}\,\delta^{(1)}_i ,
\qquad \mu_i \equiv \khat_i\cdot\nhat .
\end{equation*}

\emph{(a)} The real-space galaxy field at second order gives, as in part (i),
$b_1F_2 + \tfrac{b_2}{2} + b_{\G_2}[(\khat_1\cdot\khat_2)^2-1]$. For matter only
$F_2$ survives.

\emph{(b)} With $u_z^{(2)}(\vk) = -i\mu\,\theta^{(2)}(\vk)/k$ and
$\theta^{(2)} = -\cH f\,\Theta^{(2)}$,
\begin{equation*}
-\frac{ik\mu}{\cH}\cdot\left(-\frac{i\mu}{k}\right)\theta^{(2)}
= -\frac{\mu^2}{\cH}\,\theta^{(2)}
= f\mu^2\,\Theta^{(2)}
\;\longrightarrow\; f\mu^2 G_2(\vk_1,\vk_2) ,
\end{equation*}
using \eqref{eq:Theta_def}: it is $\Theta$, not $\theta$, whose kernel is $G_2$.

\emph{(c)} The cross term. In the convolution one leg supplies the velocity and
the other the density:
\begin{equation*}
-\frac{ik\mu}{\cH}\cdot\frac{i\cH f\mu_1}{k_1}\cdot b_1
= f k\mu\,\frac{\mu_1}{k_1}\,b_1 .
\end{equation*}
This is not symmetric in $1\leftrightarrow2$, and $Z_2$ is defined symmetrized,
so replace it by half the sum of the two orderings:
\begin{equation*}
{\rm (c)} \;\to\; \frac{f\mu k}{2}\,b_1
\left[\frac{\mu_1}{k_1} + \frac{\mu_2}{k_2}\right] .
\end{equation*}
These are exactly the $b_1$ pieces of the second line of \eqref{eq:Z2}. The
factor of $\tfrac12$ is the symmetrization, not a Taylor coefficient.

\emph{(d)} Both legs supply a velocity:
\begin{equation*}
-\frac12\left(\frac{k\mu}{\cH}\right)^{\!2}
\left(\frac{i\cH f\mu_1}{k_1}\right)\!\left(\frac{i\cH f\mu_2}{k_2}\right)
= -\frac12\,k^2\mu^2\,(i^2)\,\frac{f^2\mu_1\mu_2}{k_1k_2}
= \frac{f^2\mu^2k^2}{2}\,\frac{\mu_1\mu_2}{k_1k_2} ,
\end{equation*}
already symmetric. Here the $\tfrac12$ \emph{is} the Taylor coefficient of the
exponential, and the sign flip comes from $i^2 = -1$.

\emph{Matching.} It remains to check that (d) equals the $f^2$ pieces of
\eqref{eq:Z2}. Those pieces are
\begin{equation*}
\frac{f\mu k}{2}\left[\frac{\mu_1}{k_1}f\mu_2^2
+ \frac{\mu_2}{k_2}f\mu_1^2\right]
= \frac{f^2\mu k}{2}\,\frac{\mu_1\mu_2}{k_1k_2}
\Big[k_2\mu_2 + k_1\mu_1\Big] .
\end{equation*}
Since $\vk = \vk_1+\vk_2$, contracting with $\nhat$ gives the identity
$k\mu = k_1\mu_1 + k_2\mu_2$. The bracket is therefore $k\mu$, and the whole
expression collapses to $\tfrac{f^2\mu^2k^2}{2}\,\mu_1\mu_2/(k_1k_2)$, which is
(d). The two forms are identical, and \eqref{eq:Z2} is verified term by term.

Reading off the anatomy: $b_1F_2$ is gravity, $f\mu^2G_2$ is the linear map
acting on the second-order velocity, the $b_1$ part of the second line is
velocity $\times$ density, and the $f^2$ part is velocity $\times$ velocity. The
symmetry of the second line under $1\leftrightarrow2$ reflects that either mode
can supply the velocity.
\end{solution}

\begin{solution}{Exercise 3.3 (Kaiser multipoles).}
From \eqref{eq:kaiser},
$P^s(k,\mu) = (b_1^2 + 2b_1f\mu^2 + f^2\mu^4)\Plin(k)$. The needed moments are
\begin{equation*}
\frac12\int_{-1}^{1}\!\dd\mu\;\mu^{2n}
= \frac{1}{2n+1} ,
\end{equation*}
so $\langle\mu^0\rangle = 1$, $\langle\mu^2\rangle = 1/3$,
$\langle\mu^4\rangle = 1/5$, and for the quadrupole, using
$\mathcal{L}_2 = (3\mu^2-1)/2$,
\begin{equation*}
\frac{5}{2}\int_{-1}^{1}\!\dd\mu\,\mu^2\mathcal{L}_2 = \frac{2}{3},
\qquad
\frac{5}{2}\int_{-1}^{1}\!\dd\mu\,\mu^4\mathcal{L}_2 = \frac{4}{7}.
\end{equation*}
Hence
\begin{equation*}
P_0 = \left(b_1^2 + \frac23 b_1f + \frac15 f^2\right)\Plin ,
\qquad
P_2 = \left(\frac43 b_1f + \frac47 f^2\right)\Plin ,
\end{equation*}
and for the hexadecapole only the $\mu^4$ term survives the projection,
\begin{equation*}
P_4 = \frac{9}{2}\int_{-1}^{1}\!\dd\mu\,\mu^4\mathcal{L}_4 \cdot f^2\Plin
= \frac{8}{35}f^2\,\Plin .
\end{equation*}
Note that $P_2/P_0$ depends only on $\beta \equiv f/b_1$, which is why the
quadrupole-to-monopole ratio was the classical way to measure the growth rate.

Be careful about what follows from this. Restoring the overall amplitude, the
three expressions above depend on the cosmology only through $b_1\sigma_8$ and
$f\sigma_8$: the linear redshift-space spectrum constrains those two products
and nothing finer. It does not determine $f$, $b_1$, and the primordial
amplitude separately, and no amount of multipole information changes that,
because the degeneracy is in the amplitude rather than the angular structure.
Extracting $f$ itself requires the broadband shape to pin the cosmology, plus a
model relating $f$ to it (in \lcdm, $f \simeq \Om^{5/9}$) and an external
amplitude calibration or prior. Higher-order statistics help, but they do not
remove that dependence on assumptions.
\end{solution}

%=====================================================================
\section*{Errors worth recognising}

Not hypothetical. Each of these has been made in this document at some point,
or in the literature, and each is easier to spot than to debug.

\begin{itemize}[itemsep=3pt,topsep=2pt]

\item \textbf{Forgetting the mirror soft region in $P_{22}$.} The integrand is
symmetric under $\vq \leftrightarrow \vk-\vq$, so the soft enhancement occurs
\emph{twice}: once as $q\to0$ and once as $|\vk-\vq|\to0$. Counting it once
leaves a factor of two adrift in the infrared cancellation, and the miss is
invisible unless you check against $P_{13}$.

\item \textbf{Confusing $G_2$ with $\G_2$.} The first is the second-order
velocity kernel of \eqref{eq:F2G2}; the second is the tidal bias operator of
\eqref{eq:G2_Gamma3}. They appear on adjacent lines of $Z_2$ and are
distinguished only by typeface. In code, name them differently.

\item \textbf{Taking $\ln D_2$ when $D_2 < 0$.} With the convention used here
the second-order Lagrangian growth factor is negative, so the growth rate must
be $f_2 = \dd\ln|D_2|/\dd\ln a$. A library logarithm will return \texttt{nan};
a hand calculation will silently return a sign error.

\item \textbf{Dropping the Alcock--Paczynski volume Jacobian.} The AP mapping
rescales the amplitude as well as the wavenumbers. Because that factor is
independent of $k$ and $\mu$, it is nearly degenerate with the free amplitude,
so omitting it does not announce itself --- the damage lands on $b_1\sigma_8$
and $f\sigma_8$, and reaches $\alpha_\parallel,\alpha_\perp$ only through the
loop terms' different $\sigma_8$ scaling. That is what makes it easy to drop and
worth not dropping.

\item \textbf{Treating $1/\bar n_g$ as a prediction.} The Poisson value is a
starting guess, not a theorem: halo exclusion suppresses it and clustering
enhances it. Fixing $P_{\rm shot}$ to $1/\bar n_g$ imposes a constraint the
theory does not supply.

\item \textbf{Assuming a counterterm can absorb any missing physics.} It absorbs
only shapes lying in the retained operator basis. A good fit with a counterterm
free is not evidence that the omitted terms were harmless.

\item \textbf{Reading the residual of a fitted curve as an accuracy.} If the
counterterm was fitted to the reference you are comparing against, the agreement
in the fit range is partly guaranteed. The informative comparison is outside it.

\end{itemize}

\end{document}
